% Partie : sets-functions-relations
% Chapitre : functions
% Section : composition

\documentclass[../../../include/open-logic-section]{subfiles}

\begin{document}

\olfileid[fr]{sfr}{fun}{cmp}
\olsection{Composition des fonctions}

\begin{explain}
\oliflabeldef{sfr:fun:inv:sec}{Nous avons vu dans la \olref[inv]{sec}
que l’inverse~$f^{-1}$ d’!!a{bijection}~$f$ est lui-même une fonction.
Une autre opération sur les fonctions est la composition : on}{On}
peut définir une nouvelle fonction en composant deux fonctions $f$ et~$g$,
c’est-à-dire en appliquant d’abord $f$, puis~$g$. Cela exige que les
images et les domaines soient compatibles : l’image de~$f$ doit être
incluse dans le domaine de~$g$.
\oliflabeldef{sfr:rel:ops:sec}{Cette opération est l’analogue, pour les
fonctions, du produit relatif des relations défini dans la
\olref[rel][ops]{sec}.}{}

Un diagramme peut aider à comprendre la composition. La
\olref{fig:composition} représente deux fonctions $f\colon A \to B$
et $g\colon B \to C$, ainsi que leur composée~$(\comp{f}{g})$.
La fonction $(\comp{f}{g})\colon A \to C$ associe à chaque !!{element}
de~$A$ un !!{element} de~$C$. Voici comment déterminer lequel :
étant donnée une entrée $x \in A$, on applique d’abord $f$ à~$x$,
ce qui donne $f(x) = y \in B$, puis on applique $g$ à~$y$,
ce qui donne $g(f(x)) = g(y) = z \in C$.
\begin{figure}
  \olasset[2\olphotowidth]{assets/diagrams/composition.tikz}
  \caption{La composée $g \circ f$ de deux fonctions $f$ et~$g$.}
  \ollabel{fig:composition}
\end{figure}
\end{explain}

\begin{defn}[Composition]
Soient $f\colon A \to B$ et $g\colon B \to C$ deux fonctions.
La \emph{composée} de $f$ suivie de~$g$ est
$\comp{f}{g}\colon A \to C$, définie par
$(\comp{f}{g})(x) = g(f(x))$.
\end{defn}

\begin{ex}
Considérons les fonctions $f(x) = x + 1$ et $g(x) = 2x$.
Puisque $(\comp{f}{g})(x) = g(f(x))$, il faut, pour chaque entrée~$x$,
prendre d’abord son successeur, puis multiplier le résultat par deux.
La composée est donc donnée par $(\comp{f}{g})(x) = 2(x+1)$.
\end{ex}

\begin{prob}
Montrez que, si $f\colon A \to B$ et $g\colon B \to C$ sont toutes deux
!!{injective}s, alors $\comp{f}{g}\colon A \to C$ est !!{injective}.
\end{prob}

\begin{prob}
Montrez que, si $f\colon A \to B$ et $g\colon B \to C$ sont toutes deux
!!{surjective}s, alors $\comp{f}{g}\colon A \to C$ est !!{surjective}.
\end{prob}

\begin{prob}
Soient $f\colon A \to B$ et $g\colon B \to C$. Montrez que le graphe
de $\comp{f}{g}$ est $R_f \mid R_g$.
\end{prob}

\end{document}
